% AUTOGENERATED from hw0_math_solutions.tex by build.py -- do not edit; edit the _solutions file.
\documentclass{cs1840}
\usepackage{graphicx}

\hwnum{0}
\hwtitle{Review Problems}
\hwsubtitle{Math}
\hwdue{September 17}

\begin{document}
\makehwtitle

\hwinstructions

\begin{problem}{Policy agreement}

Read the course policies on the website.
Once finished, write ``I have read the course policies on the website'' below.
It is your responsibility to understand and follow these policies.
If further clarification is needed, please post to the discussion board.

\begin{s}
    I have read the course policies on the website    
\end{s}

\end{problem}

\begin{problem}{Bayes's rule}

\points{3} You have just finished your yearly checkup, and the doctor has good news and bad news to share.
The bad news is that you tested positive for a serious disease, and that the test is 99\% accurate
(i.e., the probability of testing positive given that you have the disease is 0.99, as is the probability of testing negative given that you don't have the disease).
The good news is that this disease is extremely rare, striking only one in 10,000 people.
What is the probability that you actually have the disease?
Show your calculations along with your final result.

\begin{s}
    Let $D$ represent having the disease and $T$ represent testing positive. By Bayes's Rule and the Law of Total Probability,

    \begin{align*}
        P(D|T)
        &= \frac{P(T|D)P(D)}{P(T)} \\
        &= \frac{P(T|D)P(D)}{P(T|D)P(D) + P(T|D^c)P(D^c)} \\
        &= \frac{0.99 \cdot 0.0001}{0.99 \cdot 0.0001 + 0.01 \cdot 0.9999} \\
        &\approx 0.0098
    \end{align*}
\end{s}

\end{problem}

\begin{problem}{Probability}

\begin{enumerate}
  \item \points{3} (independence)
    Let $X$ and $Y$ be independent random variables with PDFs $f$ and $g$, respectively.
    Consider the random variable $Z = X + Y$, and let its PDF be $h$.
    Show that
    \[ h(z) = \int_{-\infty}^\infty f(x) g(z - x) \,dx. \]
    (If you are more comfortable with discrete probabilities, you can instead derive an analogous expression for the discrete case, and give a one sentence explanation as to why your expression is analogous to the continuous case.)

  \begin{s}
    Let the CDF of $Z$ be $H(z) = P(x + y \leq z)$. Since $X$ and $Y$ are independent, their joint PDF is $f(x)g(y)$. Rearranging $x + y \leq z$ gives $y \leq z - x$. Integrating the joint PDF over this region gives

    \[
    H(z) = \int_{-\infty}^{\infty}\int_{-\infty}^{z - x} f(x)g(y) dydx = \int_{-\infty}^{\infty} f(x)G(z- x) dx
    \]

    \noindent where $G$ is the CDF of $Y$. Differentiating with respect to $z$ gives,

    \[
    h(z) = \int_{-\infty}^{\infty}f(x)g(z - x)dx
    \]
  \end{s}

  \item (conditional probabilities)
    Suppose $X$ and $Y$ are independent and uniformly distributed on $[0,1]$.
    (That is, $f(x) = g(x) = 1$ for $x \in [0,1]$ and $0$ otherwise.)
    Let $h$ be the PDF of the random variable $Z = X + Y$.

    \begin{enumerate}
      \item \points{3} What is $h$?

        \begin{s}
        The integrand from part (a) equals 1 when $0 \leq x \leq 1$ and $0 \leq z-x \leq 1$ and 0 otherwise. Rearranging the second inequality gives $z - 1 \leq x \leq z$, and taking the intersection of both inequalities gives an interval of $\max(0, z - 1) \leq x \leq \min(1, z)$. For different cases of $z$:

        \begin{itemize}
            \item $0 \leq z \leq 1$: the interval is $[0, z]$, so $h(z) = \int_0^z 1 dx = z$.
            \item $1 \leq z \leq 2$: the interval is $[z - 1, 1]$, so $h(z) = \int_{z-1}^1 1 dx = 2 - z$.
            \item $z < 0$ or $z > 2$: the interval is empty, so $h(z) = 0$.
        \end{itemize}
        \end{s}

      \item \points{3} What is
        $\Pr(X \leq 1/2 \mid X + Y \geq 5/4)$?

        \begin{s}
            By definition of conditional probability,

            \[
            P(X \leq 1/2 | X + Y \geq 5/4) = \frac{P(X \leq 1/2, X + Y \geq 5/4)}{P(X + Y \geq 5/4)}
            \]

            \noindent The joint PDF is 1 within the unit square and 0 otherwise. For the denominator, we need $x + y \geq 5/4 \Rightarrow y \geq 5/4 - x$ and $y \leq 1$, which requires $x \geq 1/4$. Integrating gives

            \begin{align*}
            P(X \leq 1/2, X + Y \geq 5/4) &= \int_{1/4}^1 \int_{5/4 - x}^1 1 dydx \\
            &= \int_{1/4}^1 \left( x - \frac{1}{4} \right) dx \\
            &= \left[ \frac{x^2}{2} - \frac{1}{4}x \right] _{1/4}^1 \\
            &= \frac{9}{32}
            \end{align*}

            \noindent For the numerator, we keep the same bounds on $y$ but also need $x \leq 1/2$. The probability equals

            \begin{align*}
            P(X \leq 1/2, X + Y \geq 5/4) &= \int_{1/4}^{1/2} \int_{5/4 - x}^1 1 dydx \\
            &= \int_{1/4}^{1/2} \left( x - \frac{1}{4} \right) dx \\
            &= \left[ \frac{x^2}{2} - \frac{1}{4}x \right] _{1/4}^{1/2} \\
            &= \frac{1}{32}
            \end{align*}

            \noindent Thus, we have

            \[P(X \leq 1/2 | X + Y \geq 5/4) = \frac{1}{9}\]
        \end{s}

    \end{enumerate}

  \item \points{3} (change of variables)
    Let $X \sim \mathcal{N}(\mu, \sigma^2)$ be Normally distributed with mean $\mu$ and variance $\sigma^2$.
    Recall that for any $a \neq 0$ and $b \in \R$, we have that $Y = aX + b$ is also Normal.
    Find $a, b$ such that $Y \sim \mathcal{N}(0,1)$.

    \begin{s}
        For $Y \sim \mathcal{N}(0,1)$, we want $E(Y) = 0$ and $\operatorname{Var}(Y) = 1$. By linearity of expectation and properties of variance, we have
    
        \begin{align*}
            E(aX + b) = 0 &\Rightarrow a\mu + b = 0 \\
            \operatorname{Var}(aX + b) = 1 &\Rightarrow a^2\sigma^2 = 1
        \end{align*}

        \noindent Solving for $a, b$ gives $a = \frac{1}{\sigma}$ and $b = -\frac{\mu}{\sigma}$ (alternatively, $a = -\frac{1}{\sigma}$ and $b = \frac{\mu}{\sigma}$ also works).
    \end{s}

  \item (covariance)
    For any two random variables $X$ and $Y$, their \textit{covariance} is defined as
    \[ \operatorname{Cov}(X,Y) := \E[(X - \E[X])(Y - \E[Y])]. \]
    (You may assume $X$ and $Y$ take on discrete values if you find it easier to work with.)

    \begin{enumerate}
      \item \points{3} Show that if $\E[Y \mid X=x] = x$, then $\operatorname{Cov}(X,Y) = \E[(X - \E[X])^2]$.

        \begin{s}
            The definition of covariance can be rearranged as
            
            \[ \operatorname{Cov}(X, Y) = E(XY) - E(X)E(Y)\]

            \noindent By Adam's law, the first term on the right equals

            \begin{align*}
                E(XY) 
                &= E(E(XY|X)) \\
                &= \sum_x E(xY|X=x)P(X=x) \\
                &= \sum_x xE(Y|X=x)P(X=x) \\
                &= \sum_x x^2P(X=x) \\
                &= E(X^2)
            \end{align*}

            \noindent and the second term on the right equals

            \begin{align*}
                E(X)E(Y) 
                &= E(X) E(E(Y|X)) \\
                &= E(X)\sum_x E(Y|X=x)P(X=x) \\
                &= E(X)\sum_x xP(X=x) \\
                &= E(X)^2
            \end{align*}

            \noindent Thus, we have $\operatorname{Cov}(X, Y) = E(X^2) - E(X)^2 = E[(X-E[X])^2]$ (where the last equality holds since $\operatorname{Var}(X) = E(X^2) - E(X)^2 = E[(X-E[X])^2]$).
        \end{s}

      \item \points{3} Show that if $X, Y$ are independent, then $\operatorname{Cov}(X,Y) = 0$.

        \begin{s}
            If $X, Y$ are independent, then

            \[
            E(XY) = \sum_x\sum_y xy P(X=x, Y=y) = \sum_xxP(X=x) \sum_yyP(Y=y) = E(X)E(Y)
            \]

            \noindent so $\operatorname{Cov}(X,Y) = E(XY) - E(X)E(Y) = 0$.
        \end{s}

    \end{enumerate}

  \item (empirical CDF)
    Let $X$ be a random variable, and let $F(x)$ be its CDF.
    Given a boolean event $A$, we can define the \textit{indicator} random variable $\1\{ A \}$ that equals $1$ if $A$ occurs, and $0$ if $A$ does not occur.
    For example, $\1\{ x \leq a \}$ is equal to $1$ when $x \leq a$, and $0$ when $x > a$.
    It follows that
    \[ F(x) = \E[\1\{X \leq x\}]. \]
    Now, let $X_1, \ldots, X_n$ be independent and identically distributed (i.i.d.) with CDF $F(x)$.
    Define \[ \widehat{F}_n(x) = \frac{1}{n} \sum_{i=1}^n \1\{X_i \leq x\}. \]

    \begin{enumerate}
      \item \points{3} For any $x \in \R$, what is $\E[ \widehat{F}_n(x) ]$?

        \begin{s}
            By linearity of expectation, we have

            \begin{align*}
                E[\widehat{F}_n(x)]
                &= E \left[\frac{1}{n} \sum_{i=1}^n \1\{X_i \leq x\} \right] \\
                &= \frac{1}{n} \sum_{i=1}^n E[\1 \{X_i \leq x \}] \\
                &= \frac{1}{n} \sum_{i=1}^n F(x) \\
                &= F(x)
            \end{align*}
        \end{s}

      \item \points{3} Show that for any $x \in \R$, we have
        \[ \Var(\widehat{F}_1(x)) = F(x)(1-F(x)). \]

        \begin{s}
            Let $I$ be an indicator random variable for $\1 \{X_1 \leq x\}$, so $E(I) = F(x)$. Since $I \in \{0, 1\}$, $I^2 = I$. So, 

            \begin{align*}
                \operatorname{Var}( \widehat{F}_1(x)) 
                &= E(\widehat{F}_1(x)^2) - E(\widehat{F}_1(x))^2 \\
                &= E(I^2) - E(I)^2 \\
                &= E(I) - E(I)^2 \\
                &= E(I) (1 - E(I)) \\
                &= F(x)(1 - F(x))
            \end{align*}
        \end{s}

      \item \points{6} Show that for any $x \in \R$, we have
        \[ \Var(\widehat{F}_n(x)) = \frac{F(x)(1-F(x))}{n}. \]

        \begin{s}
            Using the same approach and results from part (ii), we have

            \begin{align*}
                \operatorname{Var}(\widehat{F}_n{(x))}
                &= \operatorname{Var} \left(\frac{1}{n} \sum_{i=1}^n \1(X_i \leq x) \right) \\
                &= \frac{1}{n^2} \sum_{i = 1}^n \operatorname{Var}(\1 (X_i \leq x)) \\
                &= \frac{1}{n^2} \cdot n \cdot F(x)(1 - F(x)) \\
                &= \frac{F(x)(1 - F(x))}{n}
            \end{align*}
        \end{s}

      \item \points{6} Using your answer to (iii), show that
        for all $x \in \R$, we have
        \[ \E[ ( \widehat{F}_n(x) - F(x) )^2 ] \leq \frac{1}{4n}. \]

        \begin{s}
            We have

            \begin{align*}
                E[(\widehat{F}_n(x) - F(x))^2] 
                &= E[(\widehat{F}_n(x) - E[\widehat{F}_n(x)])^2] \\
                &= \operatorname{Var}(\widehat{F}_n(x)) \\
                &= \frac{F(x)(1-F(x))}{n}
            \end{align*}

            \noindent Let $p = F(x)$ and $g(p) = p(1 - p)$. Since $g(p)$ is a downwards-facing parabola with its vertex at $p = 1/2$, the maximum value of $g(p) = 1/4$. Thus, $F(x)(1 - F(x)) \leq 1/4$, so

            \[
            E[(\widehat{F}_n(x) - F(x))^2] = \frac{F(x)(1-F(x))}{n} \leq \frac{1}{4n}
            \]
        \end{s}

    \end{enumerate}

\end{enumerate}

\end{problem}

\begin{problem}{Geometry and linear algebra}

\begin{enumerate}
  \item (hyperplanes)
    A hyperplane in $\R^n$ is a set of the form $\{x \in \R^n \mid w^\top x + b = 0\}$, where $w$ is a nonzero $n$-dimensional vector and $b$ is a scalar.

    \begin{enumerate}
      \item \points{3}
        Draw the hyperplane given by $w = [-1, 2]^\top$ and $b = 2$.
        Label your axes.

        \begin{s}
            The hyperplane equation is $-x_1 + 2x_2 + 2 = 0$.

            \begin{center}
                \includegraphics[width=0.3\textwidth]{4i.png}
            \end{center}
        \end{s}

      \item \points{3}
        Draw the hyperplane given by $w = [1, 1, 1]^\top$ and $b = 0$.
        Label your axes.

        \begin{s}
            The hyperplane equation is $x_1 + x_2 + x_3 = 0$.

            \begin{center}
                \includegraphics[width=0.3\textwidth]{4ii.png}
            \end{center}
        \end{s}

      \item \points{6}
        Given a point $x_0 \in \R^n$, find the \emph{squared} distance to the hyperplane defined by $w^\top x + b = 0$.
        In other words, solve the following optimization problem.
        \begin{align*}
          \min_x \quad & \|x_0 - x\|^2\\
          \text{s.t.} \quad & w^\top x + b = 0
        \end{align*}
        (Remember: we want the squared distance, not the closest $x$.)

        \begin{s}
            We want to minimize $f(x) = \| x_0 - x \|^2$ subject to $g(x) = w^Tx+b=0$. Taking gradients of both equations gives

            \begin{align*}
                \nabla f(x) &= -2 (x_0 - x) \\
                \nabla g(x) &= w
            \end{align*}

            \noindent The Lagrange condition is $\nabla f = \lambda \nabla g$, so

            \begin{align*}
                -2(x_0 - x) = \lambda w \Rightarrow x = x_0 + \frac{\lambda w}{2}
            \end{align*}

            \noindent Applying the constraint $w^Tx + b = 0$ gives

            \begin{align*}
                w^T(x_0 + \frac{\lambda w}{2}) + b = 0 \Rightarrow \frac{\lambda}{2} = \frac{-w^Tx_0 - b}{w^Tw}
            \end{align*}

            \noindent Thus, the closest point is

            \[
            x^* = x_0 - \frac{w^Tx_0 + b}{\|w\|^2}w
            \]

            \noindent and the squared distance is 

            \[
            \|x_0 - x^*\|^2 = \left( \frac{w^Tx_0 + b}{\|w\|^2}\right)^2 \|w\|^2 = \frac{(w^Tx_0 + b)^2}{\|w\|^2}
            \]
        \end{s}

    \end{enumerate}

  \item \points{6} (rank)
    Let
      $A = \begin{bmatrix}
        1 & 2 & 1\\
        1 & 0 & 3\\
        1 & 1 & 2
      \end{bmatrix}$.

      \begin{enumerate}
        \item What is the rank of $A$? How do you know?

          \begin{s}
              Rank 2. Applying the row operations

                \[
                    R_2 \leftarrow R_2 - R_1, \quad
                    R_3 \leftarrow R_3 - R_1, \quad
                    R_2 \leftarrow -\tfrac{1}{2} R_2, \quad
                    R_3 \leftarrow R_3 + R_2
                \]

        \noindent gives
                
            \[
                \begin{bmatrix}
                    1 & 2 & 1 \\
                    1 & 0 & 3 \\
                    1 & 1 & 2
                \end{bmatrix}
                \rightarrow
                \begin{bmatrix}
                    1 & 2 & 1 \\
                    0 & 1 & -1 \\
                    0 & 0 & 0
                \end{bmatrix}
            \]
          \end{s}

        \item Find a basis for the column span of $A$.

          \begin{s}
              The pivots are in columns 1 and 2, so the corresponding columns of the original matrix form a basis for the column span of $A$:

            \[
                \left\{
                    \begin{bmatrix} 1 \\ 1 \\ 1 \end{bmatrix},
                    \begin{bmatrix} 2 \\ 0 \\ 1 \end{bmatrix}
                \right\}
            \]
          \end{s}   

      \end{enumerate}

  \item \points{6} (linear equations)
    Let
      \[ A = \begin{bmatrix}
        0 & 2 & 4 \\
        2 & 4 & 2 \\
        3 & 3 & 1
      \end{bmatrix},
      \qquad
      b = \begin{bmatrix}
        -2 \\ -2 \\ -4
      \end{bmatrix},
      \qquad\text{and}\qquad
      c = \begin{bmatrix}
        1 \\ 1 \\ 1
      \end{bmatrix}. \]

    \begin{enumerate}
      \item What is $Ac$?

        \begin{s}
            Multiplying the two matrices gives

            \[
            Ac = \begin{bmatrix}
                0 & 2 & 4 \\
                2 & 4 & 2 \\
                3 & 3 & 1
            \end{bmatrix}
            \begin{bmatrix}
                1 \\ 1 \\ 1
            \end{bmatrix}
            = 
            \begin{bmatrix}
                0 + 2 + 4 \\
                2 + 4 + 2 \\
                3 + 3 + 1
            \end{bmatrix}
            =
            \begin{bmatrix}
                6 \\ 8 \\ 7
            \end{bmatrix}
            \]
        \end{s}

      \item What is the solution to the linear system $Ax = b$?

        \begin{s}
        The linear system can be rewritten as an augmented matrix, which can be  simplified into reduced row echelon form:
        
        \[
        \left[\begin{array}{ccc|c}
            0 & 2 & 4 & -2 \\
            2 & 4 & 2 & -2 \\
            3 & 3 & 1 & -4
        \end{array}\right]
        \rightarrow
        \left[\begin{array}{ccc|c}
            1 & 0 & 0 & -2 \\
            0 & 1 & 0 & 1 \\
            0 & 0 & 1 & -1
        \end{array}\right]
        \]

        \noindent Thus, the solution is

        \[
        x =
        \begin{bmatrix}
            -2 \\ 1 \\ -1
        \end{bmatrix}
        \]
        \end{s}

    \end{enumerate}

  \item \points{3} (linear algebra)
    Let $\mat{A}, \mat{B} \in \R^{n \times n}$ and $c \in \R$, and define $f: \R^n \times \R^n \to \R$ by
    \[ f(x, y) = x^\top \mat{A} x + y^\top \mat{B} x + c. \]
    Express $f(x, y)$ in terms of $c$ and the (scalar) entries $\mat{A}_{i,j}$ and $\mat{B}_{i,j}$, using appropriate summations over the indices.

    \begin{s}
        The first term on the right expands into

        \[
        x^TAx = \sum_{i=1}^n x_i(Ax)_i = \sum_{i = 1}^n x_i \sum_{j = 1}^n A_{i, j}x_j = \sum_{i = 1}^n \sum_{j=1}^n A_{i, j}x_ix_j 
        \]

        \noindent Similarly, the second term equals

        \[
        y^TBx = \sum_{i = 1}^n \sum_{j = 1}^n B_{i, j} y_i x_j
        \]

        \noindent Thus, the overall function equals

        \[
        f(x, y) = \sum_{i = 1}^n \sum_{j=1}^n A_{i, j}x_ix_j + \sum_{i = 1}^n \sum_{j = 1}^n B_{i, j} y_i x_j + c
        \]
    \end{s}

\end{enumerate}

\end{problem}

\begin{problem}{Modeling with (finite-horizon) MDPs}
Reinforcement learning (RL) algorithms are designed to solve sequential decision-making problems, which are modeled via Markov decision processes (MDPs).
This problem is not for review, but we want you to start thinking about MDPs by choosing a problem of your own to model.

\begin{enumerate}
  \item \points{2} Think of a problem that RL could be used for. (The more creative, the better!)

    \begin{s}
        RL can be used for robotic manipulation tasks, such as teaching a robot hand to reorient a block in its palm.
    \end{s}

  \item \points{2} What are the states?

    \begin{s}
        The states are the joint angles and velocities of the hand, along with the orientation and position of the block.
    \end{s}

  \item \points{2} What are the actions?

    \begin{s}
        The actions are the target joint angles of the hand.
    \end{s}

  \item \points{2} What determines the transition function?

    \begin{s}
        The transition function is determined by physics (or a physics engine in simulation), where we calculate the probability of ending up in the target state given the starting state and action. The transition is stochastic due to unpredictable friction, motor lag, and sensor noise.
    \end{s}

  \item \points{2} How can we define rewards for this problem in terms of your states and actions?
    And what is the finite horizon?

    \begin{s}
        The reward should be the reduction in angular distance between the current and target orientations of the block. We should also include a bonus for when the target is reached and a penalty for dropping the block. For a controller running at 20 Hz, the finite horizon should be 100 timesteps, so the total process can take 5 seconds.
    \end{s}

\end{enumerate}

\end{problem}

\end{document}
